\section{Introduction}\label{sec:introduction}

An anticanonical nonvanishing theorem must turn positivity into a section of a specified line bundle. For a singular pair $(W,D)$, the line in question is the anti-log-canonical divisor
\[
P=-(K_W+D).
\]
Nefness says that $P$ has nonnegative degree on every curve. It does not itself construct an effective divisor linearly equivalent to a positive multiple of $P$. We prove such a conclusion when the pullback of $P$ has a semipositive metric with locally bounded weights and a resolution has nonzero structure-sheaf Euler characteristic. We then construct these hypotheses on a projective torus quotient. The quotient section returns upstairs as a naturally invariant anticanonical section.

All varieties are integral and defined over $\CC$. A pair $(W,D)$ is \emph{Kawamata log terminal}, abbreviated klt, if $W$ is normal, $K_W+D$ is $\QQ$-Cartier, and on a log resolution $\rho:\widetilde W\to W$, using compatible canonical divisors, every coefficient of $K_{\widetilde W}-\rho^*(K_W+D)$ is greater than $-1$. This includes the strict transforms of the boundary components as well as exceptional divisors. A normal variety is $\QQ$-factorial if every Weil divisor has a positive Cartier multiple.

A metric on a rational line bundle means the appropriate root of a metric on one Cartier multiple. In a local frame $e$ its weight is $\varphi=-\log |e|_h^2$; semipositivity means that $\varphi$ is plurisubharmonic, abbreviated psh. Locally bounded weights are bounded above and below on relatively compact coordinate neighborhoods. They need not be continuous.

We write $\PE(V)\subset N^1(V)_{\RR}$ for the closed cone generated by effective divisor classes in the real vector space of Cartier divisors modulo numerical equivalence. A rational Cartier divisor is \emph{pseudoeffective} if its class lies in this cone.

\begin{theorem}[Bounded-metric nonvanishing]\label{thm:klt}
Let $(W,D)$ be a projective $\QQ$-factorial klt pair with effective rational boundary, and put $P=-(K_W+D)$. Suppose that $P$ is nef. Suppose that on a smooth projective resolution $\rho:\widetilde W\to W$, the pullback $\rho^*P$ has a semipositive metric with locally bounded weights and
\[
\chi(\widetilde W,\OO_{\widetilde W})\ne0.
\]
Then $H^0(W,\OO_W(mP))\ne0$ for some integer $m>0$ such that $mP$ is Cartier.
\end{theorem}

These hypotheses persist on a common higher resolution: bounded psh weights pull back, and the structure-sheaf Euler characteristic is a birational invariant of smooth projective varieties. The effective boundary and $\QQ$-factoriality enter the cotangent-subsheaf and birational numerical arguments, respectively. The Euler hypothesis supplies the first twisted differential forms. None is suppressed when we pass to the quotient application.

If an algebraic torus $T$ acts on a smooth variety $Z$, its differential gives the \emph{natural linearization} on $-K_Z=\det T_Z$. A section is invariant when it is fixed by the resulting action on global sections.

\begin{theorem}[Invariant sections from a stable quotient]\label{thm:equivariant}
Let $Z$ be a smooth rationally connected projective variety whose anticanonical bundle has a smooth semipositive Hermitian metric. For every algebraic torus $T$ acting on $Z$, some integer $m>0$ satisfies
\[
H^0(Z,-mK_Z)^T\ne0,
\]
with respect to the natural anticanonical linearization.
\end{theorem}

The auxiliary GIT linearization used to choose a quotient and the
linearization in this conclusion have different roles. We perturb the anticanonical class by a small ample class and a small character to obtain a projective stable quotient. We then descend the \emph{original} naturally linearized anticanonical bundle. A character shift of that bundle would change the conclusion; Remark~\ref{rem:character} gives an elementary example.

Together with the finite-cover structure theorem, Theorem~\ref{thm:equivariant} also gives nonvanishing for every smooth projective variety with smoothly semipositive anticanonical bundle; see Corollary~\ref{cor:smooth}. The companion article \cite{SmoothCompanion} proves this smooth result by a finite-volume argument with prescribed poles. The present paper gives an alternative route through a bounded-metric criterion on klt pairs and a projective stable quotient. From the companion we use only its independently proved finite-cover structure and finite \'etale norm statements.

\subsection{Context and predecessors}

Two forms of nonvanishing must be distinguished. A divisor is \emph{num-effective} if it is numerically equivalent to an effective real divisor. A section of a specified positive multiple requires linear, rather than only numerical, equivalence. Lazi\'c--Matsumura--Peternell--Tsakanikas--Xie prove num-effectivity of nef anti-log-canonical divisors on projective log-canonical threefolds that are $\QQ$-factorial or have rational singularities \cite[Theorem~A]{LMPTX2023}. Their Theorem~E is particularly close to our Euler hypothesis: on a projective log-canonical pair with rational singularities, nef anti-log-canonical divisor and nonzero $\chi(\OO)$, every nef Cartier divisor of numerical dimension one is num-effective. Theorem~\ref{thm:klt} instead obtains a section of a positive multiple of the anti-log-canonical divisor, in arbitrary dimension, under its bounded-metric hypothesis.

For ordinary nonvanishing, M\"uller proves the result for projective klt pairs with nef anti-log-canonical divisor when its restriction to a general fiber of the maximal rationally connected fibration is semiample, and removes that additional premise in dimension three \cite[Theorem~A and Corollary~B]{Muller2025}. His equivariant theorem gives a naturally invariant section for projective sub-log-canonical pairs with semiample anti-log-canonical divisor and a commutative linear algebraic group \cite[Theorem~C]{Muller2025}. Thus our equivariant theorem uses a different positivity hypothesis while restricting the varieties to be smooth and rationally connected and the groups to be tori. The distinction between smooth semipositivity and semiampleness already occurs on surfaces: the classification of Chen--Filip--Sun--Tosatti--Zhang includes rational surfaces with smoothly semipositive anticanonical bundle and a smooth elliptic anticanonical curve whose normal bundle is nontorsion of degree zero \cite[Theorem~1.3]{CFSTZ2026}. Such an anticanonical bundle cannot be semiample, since its restriction to that curve has no nonzero section in any positive power.

The method of producing twisted differential forms from Euler characteristics and hard Lefschetz has an earlier antecedent in Demailly--Peternell--Schneider \cite[Theorem~2.7.3]{DPS2001}, there for pseudoeffective canonical bundles with algebraic-singularity metrics. Their Proposition~2.7.4 also passes from twisted sections to a determinant line. We use the hard Lefschetz theorem with multiplier ideals \cite[Theorem~2.1.1]{DPS2001}, together with the determinant argument developed by Lazi\'c--Peternell \cite[Lemma~4.1]{LP2018} and adapted to anticanonical geometry in \cite[Lemma~5.1]{LMPTX2023}. Our twisting exponents are already positive and unbounded, which permits the determinant step even when the anticanonical class is numerically trivial. To control that determinant, we use \cite[Theorem~4.1]{LMPTX2023}: for a projective $\QQ$-factorial log-canonical pair with effective boundary and nef anti-log-canonical divisor, the negative first Chern class of a saturated cotangent-tensor subsheaf is pseudoeffective. This is the log-pair extension of Ou's movable-slope method \cite{Ou2023}.

The next stage transfers metric positivity to the base of a fibration. Its inputs are the curvature theory of adjoint direct images of Berndtsson \cite{Berndtsson2009}, Berndtsson--P\u{a}un \cite{BerndtssonPaun2008}, and P\u{a}un--Takayama \cite{PaunTakayama2018}. We need the singular version because discrepancy rounding leaves fractional divisorial poles even when the original metric is bounded. Normalized fiber integrals then yield a bounded metric on a base line bundle. Skoda's small-exponent integrability theorem \cite{Skoda1972} and Fujino's Koll\'ar--Nadel vanishing theorem \cite[Theorem~1.3]{Fujino2018} allow one metric construction to control every nonnegative twist, including zero.

For the invariant theorem, projective geometric invariant theory \cite{MFK1994}, torus linearization \cite{Sumihiro1974,KKLV1989}, the attracting-cell geometry of Bia\l{}ynicki-Birula \cite{BB1973}, and Luna's slice theorem \cite{Luna1973} supply the quotient and its stabilizer boundary. M\"uller's use of fixed-point weights, convex geometry and ample perturbations is a close methodological predecessor \cite{Muller2025}. The additional analytic task is to bound the norms of translated lifts uniformly. Once that bound is proved, Kiselman's minimum principle \cite{Kiselman1978,DWZZ2018} gives the quotient metric, and the same bound permits extension of the resulting invariant section over the complement of the stable locus.

The smooth application belongs to the structure theory of manifolds with semipositive Ricci curvature. Yau's prescribed-Ricci theorem \cite{Yau1978} realizes a smooth semipositive anticanonical curvature form as a Ricci form. Demailly--Peternell--Schneider describe the resulting universal-cover decomposition and finite-cover Albanese fibration \cite[Structure Theorem, p.~218]{DPS1996}; Campana--Demailly--Peternell identify the remaining compact factors as rationally connected \cite[Theorem~1.4]{CDP2015}. The companion's finite-cover theorem retains the deck action in this decomposition. The need to descend through that action explains why ordinary nonvanishing on a rationally connected factor is insufficient and why the natural invariant conclusion of Theorem~\ref{thm:equivariant} is the relevant input.

\subsection{How the proofs fit together}

For Theorem~\ref{thm:klt}, Riemann--Roch and hard Lefschetz first produce twisted differentials at unbounded exponents. Their determinants give effective divisors
\[
N_i\sim_{\QQ}M+m_iP,\qquad m_i\longrightarrow\infty,
\qquad -M\text{ pseudoeffective}.
\]
We seek a rational fibration represented by maps $\pi:S\to W$ and $f:S\to Y$, with $S,Y$ smooth projective, $\pi$ birational and $f$ surjective with connected fibers, and a decomposition
\[
r\pi^*P=f^*V+H,
\]
where $r$ is a positive integer, $V$ is a line bundle on $Y$, and $H$ is effective. We require more than effectivity: above every prime divisor of $Y$, some component must be absent from $H$. That component will provide a nonvanishing local section with which to bound the metric transferred to $V$ near the boundary of the base.

To obtain this decomposition, choose the largest possible $\dim Y$ subject to $d\pi^*P-f^*A\in\PE(S)$ for some integer $d>0$ and ample line bundle $A$ on $Y$. Any further pencil controlled by $P$ and nonconstant on the general fiber would contradict this maximality. A prime divisor is \emph{horizontal} if it dominates $Y$, and \emph{vertical} otherwise. On a normal equidimensional model, the absence of such pencils forces the horizontal parts of the $N_i$ to vary affinely with $m_i$. Taking a difference and subtracting the base contribution determined by minimum coefficient-to-multiplicity ratios produces $V$ and $H$ with the required missing-component property.

The same maximality gives rank one for the adjoint direct images obtained by rounding the klt discrepancies. Their normalized integral metrics converge to a fiberwise \emph{essential supremum}, the appropriate notion for bounded measurable weights. The distinguished components of the preceding decomposition bound this metric near every base divisor, giving a bounded semipositive metric on $V$. The domination condition supplies a small singular perturbation with strict positivity in base directions and a controlled multiplier ideal. Vanishing for every nonnegative twist then gives an Euler polynomial on the base whose value at zero is positive: the rounded exceptional divisor has a section, and all higher cohomology vanishes. Thus a positive twist has a section. Multiplying by a suitable power of the section of $H$ and pushing back to $W$ gives a section of a positive multiple of $P$; the exceptional divisor disappears under pushforward. The initial Euler characteristic needed only to be nonzero.

For Theorem~\ref{thm:equivariant}, an ample perturbation gives a nonempty stable locus $U$ equal to the semistable locus. Its quotient $W=U/T$ is projective. Finite stabilizer slices give a klt boundary with coefficients $1-1/e$, where $e$ is the codimension-one ramification index, and identify $-K_Z|_U$ naturally with the pullback of $-(K_W+D)$. Convexity along algebraic arcs proves that translated lifts have no poles; compactification of their graph upgrades this to a locally uniform bound. The orbit supremum of the original norm consequently defines the bounded quotient metric. Rational connectedness supplies Euler characteristic one on a resolution. Theorem~\ref{thm:klt} produces a section. The bounded metric and compactness of $W$ bound that section globally. By the definition of the orbit-supremum norm, its pullback is bounded throughout $U$ and therefore extends across every omitted divisor in $Z$. No orbit estimate at an unstable point is required.

Section~\ref{sec:criterion} carries out the singular nonvanishing argument, and Section~\ref{sec:quotient} constructs the quotient and proves the invariant theorem. Section~\ref{sec:descent} gives the short smooth projective application using the finite-cover structure and norm results of the companion article \cite{SmoothCompanion}.
